\section{Límites: lo que la campaña no va a poder afirmar}
\label{sec:limites}

Esta sección existe para que ninguna de estas cosas se descubra al final. Cada una
es una afirmación que el diseño \emph{no} soporta, con la razón.

\subsection{Sobre comparaciones externas}

\textbf{Ninguna comparación con las cifras publicadas previamente por el proyecto.}
Son cantidades distintas en tres sentidos a la vez: en espacio de métrica, en
ventana temporal y en arnés de ejecución. Varias de ellas son promedios sobre los
nueve paneles o sobre categorías combinadas, que esta campaña no calcula por
construcción. Ponerlas en la misma tabla sería un error de categoría, no una
comparación desfavorable.

\textbf{Ninguna afirmación frente a métodos basados en alineamiento.} No existe hoy
un comparador de esa clase dentro de la plataforma: el único código que evalúa una
herramienta externa no está registrado como operación y su salida no puede llegar a
la tabla de resultados. La superficie de afirmación viva es \emph{interna y del lado
de la recuperación}: puede ordenar configuraciones entre sí y todavía no puede
situar ninguna frente a un método externo. Eso es exactamente lo que el despliegue
en LAFA resuelve, y es la razón por la que el despliegue es parte del diseño y no un
apéndice.

\subsection{Sobre los nodos bloqueados}

Seis de los diez nodos están bloqueados, cada uno por un cero o por un uno con
nombre, y los ceros no son uno por nodo: no hay
anotaciones de dominios, no hay mapeo de dominios a términos, no hay configuraciones
de puntuación registradas, no hay modelos de reordenación, no hay conjuntos de
datos exportados y no hay coocurrencia de términos. Que estén bloqueados no es un
fallo de la campaña: es el estado del registro, y avanzar uno por uno es el plan.

Lo que sí hay que corregir es que cuatro de esos seis bloqueos son literales en el
código y no lecturas del registro, de modo que imprimirían la misma palabra con un
millón de filas presentes. Un bloqueo estructural no debe poder leerse como
evidencia. Es el cambio 10 de la sección~\ref{sec:procedimiento}.

\subsection{Sobre reproducibilidad}

Lo medido es que la \emph{puntuación} es determinista: 117 de 117 métricas se
reproducen exactamente. Eso es más débil que decir que una corrida reproduce
modelos idénticos, porque hoy no hay modelos. La afirmación correcta es sobre el
scoring, y no debe ampliarse.

\subsection{Sobre lo que no se ha establecido}

\begin{itemize}[leftmargin=1.4em,itemsep=0.3em]
  \item \textbf{Las dos variantes de la ventana difieren} en un factor de 2.45 en
        uno de los paneles y son idénticas en otro. El mecanismo \emph{no está
        establecido}. Cualquier veredicto que las mezcle es nulo, y cualquiera que
        se publique sobre una tiene que nombrar cuál.
  \item \textbf{La ventana retira 426{,}385 anotaciones sobre 47{,}455 proteínas},
        el doble de proteínas de las que ganan anotación, y nadie ha desglosado esa
        cifra por causa: término obsoleto, cambio de relación, o retracción real.
        Hasta que se desglose, la confianza en la ventana está acotada por un número
        sin examinar mayor que las ganancias que la campaña mide.
  \item \textbf{Tres configuraciones no son evaluables}: tres capas intermedias de
        uno de los modelos producen puntuación constante. No deben entrar en ninguna
        rejilla.
  \item \textbf{La fracción de columnas vacías del vector de rasgos no está
        establecida con precisión.} Dos muestreos sucesivos discrepan. Importa para
        la ablación: retirar una columna que es nula en la mayoría de las filas
        devuelve \emph{no contribuye} por una razón que no tiene nada que ver con la
        señal.
\end{itemize}

\subsection{Y una propiedad del diseño que conviene no perder}

El holdout se protege por ausencia y no por disciplina: las releases que no deben
mirarse no están cargadas. Cualquier decisión futura que cargue una de ellas por
comodidad convierte una garantía estructural en una promesa, y las promesas de esta
clase ya se han incumplido en este proyecto sin que nadie lo notase. La fricción es
la característica, no el defecto.
